% s01-intro.tex —— 第 1 章 引言（底本印刷页 1–15）

\chapter{引言}
\label{chap:intro}

\section{背景}
\label{sec:1.1}

\subsection{紧生成}
\label{sec:1.1.1}
\srcpage{1}

[DG15] 中，Drinfeld 与 Gaitsgory 证明：DG 范畴 $\Dmod(\Bun_G)$ 具有一种名为紧生成的小性。
粗略地说，范畴 $\mathbf{C}$ 的紧对象，是那些「从 $c$ 出发的映射行为良好」的对象
$c\in\mathbf{C}$，即函子 $\Maps_{\mathbf{C}}(c,-)$ 与余极限交换；
若 $\mathbf{C}$ 的每个对象都能写成紧对象的余极限，便称 $\mathbf{C}$ 是紧生成的。%
\yizhu{「紧对象」在本库既有译本中统一作「紧对象」；「紧生成」即每个对象都是紧对象的余极限。
DG 范畴的紧对象通常也叫有限性对象。}{I：本库术语表（\texttt{compact object} $=$ 紧对象）}

紧生成是 DG 范畴上极受青睐的有限性条件，理由有二：一在范畴内部，一在范畴外部。

第一条理由是：若 $\mathbf{C}$ 紧生成，则要理解任意对象之间的映射，只需理解紧对象之间的映射。
具体地说，设 $c,d\in\mathbf{C}$，可写成
\[
c=\colim_{i\in I}c_i,\qquad d=\colim_{j\in J}d_j,
\]
其中 $c_i$ 与 $d_j$ 都是紧对象。于是有
\srcpage{2}
\begin{align}
\Maps_{\mathbf{C}}(c,d)
&=\Maps_{\mathbf{C}}\Bigl(\colim_{i\in I}c_i,\ \colim_{j\in J}d_j\Bigr)\\
&=\lim_{i\in I}\Maps_{\mathbf{C}}\Bigl(c_i,\ \colim_{j\in J}d_j\Bigr)\\
&=\lim_{i\in I}\colim_{j\in J}\Maps_{\mathbf{C}}(c_i,d_j).
\end{align}

第二条理由是：若 $\mathbf{C}$ 紧生成，则从 $\mathbf{C}$ 出发的（连续）函子，由它在紧对象上的取值决定。
换言之，若 $\mathbf{D}$ 是另一个（不必紧生成的）余完备 DG 范畴，则有
\[
\Fun_{\mathrm{cts}}(\mathbf{C},\mathbf{D})=\Fun_{\mathrm{exact}}(\mathbf{C}^{c},\mathbf{D}),
\]
其中 $\mathbf{C}^{c}$ 是紧对象构成的全子范畴。

以下是一个动机，说明我们为何特别想研究 $\Dmod(\Bun_G)$。
（de Rham、整体、临界层级的）非分歧几何 Langlands 猜想断言：存在 DG 范畴的等价
\[
\Dmod(\Bun_G)\xrightarrow{\ \sim\ }\IndCoh_{\Nilp}(\LocSys_{\Gd}),
\]
并满足若干相容性条件。%
\yizhu{$\IndCoh$ 保留记号，中文可作「归纳凝聚层」；下标 Nilp 指奇异支撑落在幂零锥内。}{I：本库 SY02-29 处理方式}
等价右端的 $\IndCoh_{\Nilp}(\LocSys_{\Gd})$ 等同于 ind-完备化 $\Ind(\Coh_{\Nilp}(\LocSys_{\Gd}))$，
故先验是紧生成的。于是，$\Dmod(\Bun_G)$ 紧生成是几何 Langlands 猜想成立的必要条件。

该猜想近来由九位作者的系列五篇文章 [GR24a]、[ABC+24a]、[CCF+24]、[ABC+24b]、[GR24b] 证明：
Arinkin、Beraldo、Campbell、Chen、Faergeman、Gaitsgory、Lin、Raskin 与 Rozenblyum。
证明的第一步见 [GR24a]，即构造 Langlands 函子
\srcpage{3}
\[
\Lang_G:\Dmod(\Bun_G)\longrightarrow\IndCoh_{\Nilp}(\LocSys_{\Gd}).
\]
紧生成在这里的用处如下：存在一个幺半作用
\[
\QCoh(\LocSys_{\Gd})\otimes\Dmod(\Bun_G)\longrightarrow\Dmod(\Bun_G),
\]
称为谱作用。固定其中一个对象
\[
\Poinc^{\mathrm{Vac},\mathrm{glob}}_{G,!}\in\Dmod(\Bun_G),
\]
即真空 Poincaré 层，便得到函子
\[
\QCoh(\LocSys_{\Gd})\longrightarrow\Dmod(\Bun_G).
\]
该函子有连续右伴随，记作
\[
\Lang_{G,\mathrm{coarse}}:\Dmod(\Bun_G)\longrightarrow\QCoh(\LocSys_{\Gd}).
\]
[GR24a, 定理 1.6.2] 证明：$\Lang_{G,\mathrm{coarse}}$ 把 $\Dmod(\Bun_G)$ 的紧对象送到
$\QCoh(\LocSys_{\Gd})$ 的最终余连通对象。换言之，有正合函子
\[
\Dmod(\Bun_G)^{c}\longrightarrow\QCoh(\LocSys_{\Gd})^{>-\infty}.
\]
又因函子
\[
\Psi^{\Nilp,\{0\}}:\IndCoh_{\Nilp}(\LocSys_{\Gd})\longrightarrow\QCoh(\LocSys_{\Gd})
\]
在最终余连通的全子范畴上是等价，遂得
\[
\Dmod(\Bun_G)^{c}\longrightarrow\IndCoh_{\Nilp}(\LocSys_{\Gd})^{>-\infty}
\subset\IndCoh_{\Nilp}(\LocSys_{\Gd}).
\]
最后，由于 $\Dmod(\Bun_G)$ 紧生成，便得到连续函子
\[
\Lang_G:\Dmod(\Bun_G)\longrightarrow\IndCoh_{\Nilp}(\LocSys_{\Gd}).
\]

一个自然的追问与推广方向，是从非分歧情形走向分歧情形。
也就是说，不再考虑叠 $\Bun_G$，而考虑 $X$ 上带某点 $x\in X$ 处层级结构的主 $G$-丛的模叠。
\srcpage{4}
设 $D_x$ 是 $x\in X$ 周围的形式圆盘。把主 $G$-丛 $\mathcal{F}_G$ 沿 $D_x\to X$ 限制，
得到 $D_x$ 上的主 $G$-丛 $\mathcal{F}_{G,D_x}$。
于是可把 $G(\mathcal{O})$ 解释为常值群概形 $D_x\times G$ 在 $D_x$ 上的截面群，
并把 $\mathcal{F}_{G,D_x}$ 看作 $G(\mathcal{O})$-挠子。
对子群 $H\subset G(\mathcal{O})$，$\mathcal{F}_G$ 上的（$x$ 处）$H$-层级结构，
指 $\mathcal{F}_{G,D_x}$ 到 $H$ 的一个约化。
$X$ 上带 $x\in X$ 处 $H$-层级结构的主 $G$-丛的模叠记作 $\Bun^H_G$。

顺带一提：当然可以在任意有限点集 $x_1,\dots,x_k\in X$ 处考虑 $H$-层级结构，
甚至每点取不同的子群。但为叙述简便，此处及后文只讨论单点情形，必要时再提多点情形。

不妨指出：当 $H=G(\mathcal{O})$ 时，$H$-层级结构不增添任何结构，即
$\Bun^{G(\mathcal{O})}_G=\Bun_G$。
子群 $H$ 越小，分歧越野；例如 $H=\{e\}$ 的情形称为全层级结构，
相应的叠记作 $\Bun^{\infty x}_G$。

这就引出自然的问题：对任意子群 $H\subset G(\mathcal{O})$，
能否用谱的语言描述范畴 $\Dmod(\Bun^H_G)$？也就是说，是否存在分歧的几何 Langlands 等价？

对任意子群 $H\subset G(\mathcal{O})$，以现有技术无法回答此问题。
因此作为第一步，考虑 $G(\mathcal{O})$ 的大子群（换言之，驯顺分歧）是合理的。
更确切地说，有（满的）群同态
\[
G(\mathcal{O})\longrightarrow G,
\]
由在 $t=0$ 处取值给出，其核称为第一同余子群 $G^{(1)}(\mathcal{O})$。
若子群 $H\subset G(\mathcal{O})$ 包含 $G^{(1)}(\mathcal{O})$，就称 $H$ 是驯顺的。%
\yizhu{「驯顺」对应英文 tame；此处是子群的性质，与分歧的「驯顺/野性」用词一致。}{A：底本 §1.1.1 p.4}

驯顺的几何意义在于：对驯顺子群 $H$，$H$-层级结构等同于 $x$ 处的一个
$H/G^{(1)}(\mathcal{O})$-约化。
用大白话说，驯顺层级结构只活在点本身，而不活在形式邻域上。

我们将着眼于驯顺子群 $I$，称为 Iwahori 子群。
它是 Borel 子群 $B\subset G$ 在同态 $G(\mathcal{O})\to G$ 下的原像。%
\yizhu{本库各工程对 Iwahori 一词一律音译保留，作「Iwahori 子群」。}{I：本库术语表（\texttt{Iwahori subgroup} $=$ Iwahori 子群）}
于是 Iwahori 层级结构就是 $x\in X$ 处的一个 $B$-约化。

作为走向 Iwahori 分歧几何 Langlands 猜想的一步，本论文将证明
$\Dmod(\Bun^I_G)$ 紧生成。

本文按如下方式建立在 [DG15] 之上：不再考虑 $\Bun_G$，而考虑分类 $X$ 上主 $G$-丛
连同固定点 $x\in X$ 处 Iwahori 层级结构（即固定点 $x\in X$ 处 $B$-约化）的代数叠 $\Bun^I_G$。
主要结果是
\srcpage{5}
\begin{theorem}[\textup{[Nam24]}]
\label{thm:1.1.1.1}
$\Dmod(\Bun^I_G)$ 是紧生成的。
\end{theorem}

更一般地，对有限点集 $x_1,\dots,x_k\in X$，可考虑分类 $X$ 上主 $G$-丛，
连同每个 $x_j$ 处 Iwahori 层级结构的代数叠 $\Bun^{(I;x_1,\dots,x_k)}_G$。
我们的论证在此一般性下几乎逐字适用，从而证明
\begin{theorem}
\label{thm:1.1.1.2}
$\Dmod(\Bun^{(I;x_1,\dots,x_k)}_G)$ 是紧生成的。
\end{theorem}

有人向我们指出，Gaitsgory 与 Arkhipov 的未发表文章 [AG09] 概述了非分歧与
Iwahori 分歧两种情形下紧生成的证明。然而该论证有一处组合性质的缺陷，
[DG15] 已在非分歧情形将其修正。我们的结果表明，对 Iwahori 分歧情形，
类似的修改已经足够。

\subsection{Iwahori 基本局部等价}
\label{sec:1.1.2}

（de Rham、整体、临界层级的）非分歧几何 Langlands 猜想近来由九位作者的
系列五篇文章 [GR24a]、[ABC+24a]、[CCF+24]、[ABC+24b]、[GR24b] 证明。
其陈述是：存在 DG 范畴的等价
\[
\Dmod(\Bun_G)\xrightarrow{\ \sim\ }\IndCoh_{\Nilp}(\LocSys_{\Gd}),
\]
并满足若干相容性。

在第一篇 [GR24a] 中，作者构造了 Langlands 函子
\[
\Lang_G:\Dmod(\Bun_G)\longrightarrow\IndCoh_{\Nilp}(\LocSys_{\Gd}).
\]
构造的主要构件是谱作用
\[
\QCoh(\LocSys_{\Gd})\otimes\Dmod(\Bun_G)\longrightarrow\Dmod(\Bun_G)
\]
以及 $\Dmod(\Bun_G)$ 是紧生成的 DG 范畴这一事实。

在第二篇 [ABC+24a] 中，Langlands 函子在局部得到研究。更确切地说，指的是下面这件事。
\srcpage{6}
在点 $x\in X$ 处，可考虑三个局部性质的等价。它们是几何 Casselman–Shalika 等价
\[
\CS_{G,x}:\Whit^{!}(G)_x\xrightarrow{\ \sim\ }\Rep(\Gd)_x,
\]
几何 Satake 等价
\[
\Sat_{G,x}:\Sph_{G,x}\xrightarrow{\ \sim\ }\Sph^{\spec}_{\Gd,x},
\]
以及（临界层级的）基本局部等价
\srcpage{7}
\[
\FLE_{G,x}:\KL(G)_{\mathrm{crit},x}\xrightarrow{\ \sim\ }
\IndCoh(\Op^{\monfree}_{\Gd,x}).
\]
此外，几何 Satake 等价是\emph{幺半}范畴的等价。
$\Sph_{G,x}$ 作用在 $\Whit^{!}(G)_x$（相应地 $\KL(G)_{\mathrm{crit},x}$）上，
$\Sph^{\spec}_{\Gd,x}$ 作用在 $\Rep(\Gd)_x$（相应地 $\IndCoh(\Op^{\monfree}_{\Gd,x})$）上；
在 $\Sat_{G,x}$ 之下，$\CS_{G,x}$（相应地 $\FLE_{G,x}$）与这些作用相容。

上述文献还证明，这些等价与 Langlands 函子相容，意思是有交换图
\begin{center}
\begin{tikzcd}[row sep=large,column sep=huge]
\Whit^{!}(G)_x \arrow[r,"\CS_{G,x}"] \arrow[d] & \Rep(\Gd)_x \arrow[d]\\
\Dmod(\Bun_G) \arrow[r,"\Lang_G"] \arrow[d] & \IndCoh_{\Nilp}(\LocSys_{\Gd}) \arrow[d]\\
\KL(G)_{\mathrm{crit},x} \arrow[r,"\FLE_{G,x}"] & \IndCoh(\Op^{\monfree}_{\Gd,x}).
\end{tikzcd}
\end{center}

\srcpage{8}
此外，$\Sph_{G,x}$ 作用在 $\Dmod(\Bun_G)$ 上，$\Sph^{\spec}_{\Gd,x}$ 作用在
$\IndCoh_{\Nilp}(\LocSys_{\Gd})$ 上；在 $\Sat_{G,x}$ 之下，$\Lang_G$ 与这些作用相容。

然而，这还不足以控制整体 Langlands 函子的行为。
与其选定一点 $x\in X$，不如同时考虑 $X$ 的所有点。
不仅如此，还要允许这些点在 $X$ 中「移动」，换言之，允许 $x$ 以族的形式变动。
此外，还要考虑点在 $X$ 中移动时「彼此碰撞」的情形；换言之，
允许 $X$ 中点的多元组构成的族。

事实证明，编码上述信息的正确技术装置是\emph{因子化}。
\yizhu{英文 literature 里有两个易混的概念：「factorization algebra」（nLab 词条所述，
是 Costello--Gwilliam 意义下量子场论的分解代数）与本译的「因子化范畴」
（Ran 空间上的范畴层）。本文将后者一律译作「因子化范畴」，遇到前者时随文注明。}{B：nLab, 「factorization algebra」}
上述等价都是因子化范畴的等价在点 $x\in\Ran$ 处的限制：
\[
\CS_G:\Whit^{!}(G)\xrightarrow{\ \sim\ }\Rep(\Gd),\qquad
\Sat_G:\Sph_G\xrightarrow{\ \sim\ }\Sph^{\spec}_{\Gd},
\]
以及
\[
\FLE_G:\KL(G)_{\mathrm{crit}}\xrightarrow{\ \sim\ }\IndCoh(\Op^{\monfree}_{\Gd}).
\]

若在整体几何 Langlands 的语境中，考虑固定点 $x_0\in X$ 处的驯顺（Iwahori）分歧，
自然会问：能否给出带 Iwahori 层级结构的 $\Bun_G$ 上 D-模范畴的谱描述？
换言之，$\Dmod(\Bun^I_G)\simeq{?}$ 的右端应该是什么？

我们有很强的启发式理由期待，答案应是叠
\[
\LocSys_{\Gd}(X\setminus x_0)
\times_{\LS_{\Gd}(D^\times_{x_0})}\LS_{\Bd}(D^\times_{x_0})
\times_{\LS_{\Td}(D^\times_{x_0})}\LS_{\Td}(D_{x_0})
\]
上带某种奇异支撑条件（我们仍称之为 Nilp）的归纳凝聚层。%
\yizhu{原文把 $\LS$ 与 $\LocSys$ 混用：局部系统叠记 $\LS_{\Gd}$，整体记 $\LocSys_{\Gd}$；
按底本记号照排，不统一。三轮纤维积的基底依次是 $\LS_{\Gd}(D^\times_{x_0})$、
$\LS_{\Td}(D^\times_{x_0})$。}{A：底本 §1.1.2 p.9 与 §3.1.2}

\srcpage{9}
利用 Faergeman 在 [Fae24] 中构造的驯顺分歧谱作用，以及我们在 [Nam24] 中
对 $\Dmod(\Bun^I_G)$ 紧生成的证明，可以仿照 [GR24a] 构造函子
\[
\Lang^I_G:\Dmod(\Bun^I_G)\xrightarrow{\ \sim\ }
\IndCoh_{\Nilp}\Bigl(\LocSys_{\Gd}(X\setminus x_0)
\times_{\hspace{-0.35em}\LS_{\Gd}(D^\times_{x_0})}\hspace{-0.35em}\LS_{\Bd}(D^\times_{x_0})
\times_{\hspace{-0.35em}\LS_{\Td}(D^\times_{x_0})}\hspace{-0.35em}\LS_{\Td}(D_{x_0})\Bigr).
\]

沿五篇文章的策略，下一个合乎逻辑的步骤是在局部研究 $\Lang^I_G$。
由于在 $x_0$ 之外该函子是非分歧的，更确切地说，应在 $x_0$ 处局部地研究它。
非分歧情形的三个逐点等价都有自然的 Iwahori 分歧对应物：
Arkhipov–Bezrukavnikov 等价
\[
\AB_{x_0}:\Whit^{!}(\Fl_G)_{x_0}\xrightarrow{\ \sim\ }\QCoh(\nd/\Bd)_{x_0},
\]
Bezrukavnikov 等价
\[
\B_{x_0}:\Aff_{G,x_0}\xrightarrow{\ \sim\ }\Aff^{\spec}_{\Gd,x_0},
\]
以及 Iwahori 基本局部等价
\[
\IFLE_{x_0}:\IKL(G)_{\mathrm{crit},x_0}\xrightarrow{\ \sim\ }
\IndCoh\Bigl(\Op^{\mer}_{\Gd,x_0}
\times_{\LS_{\Gd}(D^\times_{x_0})}\LS_{\Bd}(D^\times_{x_0})
\times_{\LS_{\Td}(D^\times_{x_0})}\LS_{\Td}(D_{x_0})\Bigr).
\]

我们期待这些逐点等价与 $\Lang^I_G$ 以完全类似于非分歧情形的方式相互作用，
但同样不期待这已经足够。
我们要允许 $x\in X$ 以族的形式变动，但 Iwahori 子群不可因子化，
所以不让 $x_0$ 四处移动。
实际上，我们要允许 $X$ 中点的多元组构成的族，其中包含取常值 $x_0$ 的常值族，
这样就能考虑「移动」点 $x$ 与「静止」点 $x_0$「碰撞」时会发生什么。

\srcpage{10}
此情形下我们建议采用的技术装置是\emph{因子化模范畴}理论。
更确切地说，我们将考虑\emph{成对}范畴 $(\mathcal{C},\mathcal{M})$，
其中 $\mathcal{C}$ 是因子化范畴，$\mathcal{M}$ 是 $x_0$ 处的因子化 $\mathcal{C}$-模范畴。

本文证明：$x_0$ 处六个局部 Iwahori 分歧范畴中的每一个，
都可通过与其非分歧对应物相应的因子化范畴，赋予因子化模范畴的结构。
此外，我们把三个逐点 Iwahori 分歧等价提升为成对等价的等价
\[
\AB:\bigl(\Whit^{!}(G),\Whit^{!}(\Fl_G)\bigr)\xrightarrow{\ \sim\ }
\bigl(\Rep(\Gd),\QCoh(\nd/\Bd)\bigr),
\]
\[
\B:\bigl(\Sph_G,\Aff_G\bigr)\xrightarrow{\ \sim\ }
\bigl(\Sph^{\spec}_{\Gd},\Aff^{\spec}_{\Gd}\bigr),
\]
以及
\begin{multline*}
\IFLE:\bigl(\KL(G)_{\mathrm{crit}},\IKL(G)_{\mathrm{crit}}\bigr)\xrightarrow{\ \sim\ }\\
\Bigl(\IndCoh^{*}(\Op^{\monfree}_{\Gd}),\
\IndCoh^{*}\bigl(\Op^{\mer}_{\Gd}\times_{\LS^{\mer}_{\Gd}}\nd/\Bd\bigr)\Bigr).
\end{multline*}

\srcpage{11}
在 §3.1 中，我们构造所关心的六个因子化模范畴。
我们把它们看成并记作成对范畴 $(\mathcal{C},\mathcal{M})$，
其中 $\mathcal{C}$ 是因子化范畴，$\mathcal{M}$ 是（$x_0$ 处的）因子化 $\mathcal{C}$-模范畴。
在全部六种情形，因子化范畴（即第一分量）都在 [ABC+24a] 中定义。

\begin{proposition}[\textup{[ABC+24a]}]
\label{prop:1.1.2.1}
存在因子化范畴的等价
\[
\CS_G:\Whit^{!}(G)\xrightarrow{\ \sim\ }\Rep(\Gd),
\]
称为几何 Casselman–Shalika 等价，且落在 $\FactAlgCat$ 中。
\end{proposition}

在 §3.2 中，我们构造成对范畴的态射
\[
\AB:\bigl(\Whit^{!}(G),\Whit^{!}(\Fl_G)\bigr)\longrightarrow
\bigl(\Rep(\Gd),\QCoh(\nd/\Bd)\bigr),
\]
其第一分量是 $\CS_G$。我们还证明

\begin{theorem}
\label{thm:1.1.2.2}\label{thm:AB}
成对范畴的态射 $\AB$ 是 $\FactModCat_{x_0}$ 中的等价。
\end{theorem}

\begin{proposition}[\textup{[ABC+24a]}]
\label{prop:1.1.2.3}
存在幺半等价
\[
\Sat_G:\Sph_G\xrightarrow{\ \sim\ }\Sph^{\spec}_{\Gd},
\]
称为（导出）几何 Satake 等价，且落在 $\mathrm{Alg}(\FactAlgCat)$ 中。
等价 $\CS_G$ 关于 $\Sat_G$ 等变。
\end{proposition}

在 §3.3 中，我们构造幺半成对范畴的态射
\[
\B:\bigl(\Sph_G,\Aff_G\bigr)\longrightarrow
\bigl(\Sph^{\spec}_{\Gd},\Aff^{\spec}_{\Gd}\bigr),
\]
其第一分量是 $\Sat_G$。我们还证明

\begin{theorem}
\label{thm:1.1.2.4}\label{thm:B}
幺半成对范畴的态射 $\B$ 是 $\mathrm{Alg}(\FactModCat_{x_0})$ 中的等价，
且成对等价 $\AB$ 关于它等变。
\end{theorem}

\begin{proposition}[\textup{[ABC+24a]}]
\label{prop:1.1.2.5}
存在因子化范畴的等价
\[
\FLE:\KL(G)_{\mathrm{crit}}\xrightarrow{\ \sim\ }\IndCoh^{*}(\Op^{\monfree}_{\Gd}),
\]
称为基本局部等价，且关于 $\Sat_G$ 等变。
\end{proposition}

在 §3.4 中，我们构造成对范畴的态射
\[
\IFLE:\bigl(\KL(G)_{\mathrm{crit}},\IKL(G)_{\mathrm{crit}}\bigr)\longrightarrow
\Bigl(\IndCoh^{*}(\Op^{\monfree}_{\Gd}),\
\IndCoh^{*}\bigl(\Op^{\mer}_{\Gd}\times_{\LS^{\mer}_{\Gd}}\nd/\Bd\bigr)\Bigr),
\]
其第一分量是 $\FLE$。我们还证明

\srcpage{12}
\begin{theorem}
\label{thm:1.1.2.6}
成对范畴的态射 $\IFLE$ 是 $\FactModCat_{x_0}$ 中的等价，且关于 $\B$ 等变。
\end{theorem}

\section{论文的结构}
\label{sec:1.2}

本论文是装订成册的学位论文，其组成部分是 [Nam24] 与 [Nam26]。
第 2 章是 [Nam24] 的主体，第 3 章是 [Nam26] 的主体。
两篇文章的引言经过编辑合并，成为本论文的引言。

\subsection{第 2 章的结构}
\label{sec:1.2.1}

下面说明第 2 章的组织方式。

在 §2.1 中，我们首先在 §2.1.1 给出紧生成 DG 范畴的定义。
然后概述 Drinfeld 与 Gaitsgory 在 [DG15] 中对 $\Dmod(\Bun_G)$ 紧生成的证明。
论证的极简版本如下。

首先，同一批作者在 [DG13] 中的先前结果表明：QCA 代数叠，
即拟紧（「QC」）且自同构群仿射（「A」）的代数叠，
其 D-模范畴是紧生成的。我们在 §2.1.2 概述这一讨论。

麻烦在于 $\Bun_G$ 不是 QCA 的，因为它不拟紧。
不过 $\Bun_G$ 的自同构群确实仿射，所以它是局部 QCA 的。
这意味着 $\Bun_G$ 是一些拟紧 QCA 开子叠的并，故 $\Dmod(\Bun_G)$ 是紧生成 DG 范畴
沿 $!$-拉回构成的极限。

若这些 $!$-拉回都有定义良好的左伴随，则该极限等同于以 $!$-拉回的左伴随
（即 $!$-推前函子）为连接态射的余极限。
于是，由于这些 $!$-推前函子有连续右伴随，该余极限会因形式理由自动紧生成。
不幸的是，$!$-推前函子并非沿所有映射都有定义。

开子叠的一个称为\emph{余截断性}的性质保证：沿包含映射 $U\to U'$ 的
$!$-推前函子都有定义。这在 §2.1.3 讨论。

随后我们在 §2.1.4 回顾 $U$ 的闭补上的一个几何条件，称为\emph{收缩性}，
它足以保证 $U$ 余截断。

[DG15] 给出的 $\Bun_G$ 到拟紧开子叠的明确分层，粗略地说称为 Harder–Narasimhan 分层，
与主 $G$-丛到抛物子群的约化的次数有关。我们在 §2.1.5 描述这一分层。

最后，我们在 §2.1.6 概述：上述拟紧开子叠的闭补确实收缩
（从而拟紧开子叠确实余截断）。

在 §2.2 中，我们对 $G=\mathrm{SL}_2$ 走一遍紧生成的证明，
既在非分歧情形（§2.2.1），也在 Iwahori 分歧情形（§2.2.2）。
证明的大部分概念性想法已在此玩具例中出现，因此，想对证明如何运作有大致了解、
又不想第一遍就细究组合学的读者，可以跳过 §2.1.5 与 §2.1.6，直接读 §2.2。

在 §2.3 中，我们证明本文的主要结果，即 \thmref{thm:1.1.1.1}：
$\Dmod(\Bun^I_G)$ 紧生成。

我们用 §2.3.1 与 §2.3.2 给出 $\Bun^I_G$ 的类似分层。
粗略地说，先按底主 $G$-丛在 Harder–Narasimhan 分层中的层来分层，
再按 Iwahori 层级结构（即点处的 $B$-约化）与 $P$-约化的相对位置进一步分层。

随后我们用 §2.3.3 至 §2.3.7 证明：此分层中各拟紧开子叠的闭补确实收缩。
最后，我们在 §2.3.8 讨论 \thmref{thm:1.1.1.2}，即
$\Dmod(\Bun^{(I,x_1,\dots,x_k)}_G)$ 紧生成的证明；
说明整个论证只有一处与分层定义中某个数值界有关的改动，其余照搬即可。

\subsection{第 3 章的结构}
\label{sec:1.2.2}

下面说明第 3 章的组织方式。

我们首先在 §3.1.1 回顾因子化的基础。这样做主要是为了固定记号；
细节请读者参阅 [Ras]。

在 §3.1.2、§3.1.3 与 §3.1.4 中，我们分别定义出现在 $\AB$、$\B$ 与 $\IFLE$
两侧的六个因子化模范畴。

在 §3.1.5 中，我们证明 \thmref{thm:3.1.5.2}，
它使我们可以用因子化模的范畴来逼近各种因子化模范畴。
这一点将用于构造 $\B$ 与 $\IFLE$。
它是 \propref{prop:3.1.5.1} 的因子化模范畴版本：后者粗略地说，
用因子化模的范畴逼近因子化范畴，已在 [ABC+24a] 中证明。

接下来的三节分别构造等价 $\AB$、$\B$ 与 $\IFLE$。

在 §3.2.1 中，我们讨论\emph{可熔性}这一概念，它将出现在 Bogdanova、Chen 与 Gaitsgory
即将发表的文章 [BCG26] 中。
可熔性是关于因子化 $\Rep(\Gd)$-模范畴的一个可组合地检验的条件，
它使我们可以把「因子化 $\Rep(\Gd)$-模范畴之间的函子是等价」化归为
「$x_0$ 处的纤维是普通范畴的等价」。我们论证 $\Whit^{!}(\Fl_G)$ 可熔。

在 §3.2.2 中，我们利用 [CCF+24] 中构造的半无限版本，构造成对范畴的态射 $\AB$。
上述文献还证明 $\AB$ 的半无限版本是逐点等价，
这使我们能在 §3.2.3 断定 $\AB_{x_0}$ 是等价，从而证明 \thmref{thm:1.1.2.2}。

在 §3.3.1 中，我们证明：沿 $\Rep(\Gd)\to\Sph^{\spec}_{\Gd}$ 的因子化限制 2-函子
（定义在 Chen 与 Fu 的未发表文档 [CF] 中）是保守的。
由于 $\Aff_G$ 在此函子下的像可熔，这使我们可以把「$\B$ 是等价」
化归为「$\B_{x_0}$ 是等价」。

在 §3.3.2 中，我们利用 $\Aff_G$ 经 $\AB$ 对 $\QCoh(\nd/\Bd)$ 的作用，
构造幺半成对范畴的态射 $\B$。
为证明 $\B_{x_0}$ 是等价，我们在 §3.3.3 把 Dhillon 与 Taylor 在 [DT25] 中的想法
从 Betti 情形改造过来。这就证明了 \thmref{thm:1.1.2.4}。

在 §3.4.1 中，我们构造 $\IFLE$。
在 §3.4.2 中，我们论证 $\IKL(G)_{\mathrm{crit}}$ 紧生成、且 $\IFLE$ 保持紧对象，
以便把「$\IFLE$ 是等价」化归为「$\IFLE_{x_0}$ 是等价」。

在 §3.4.3 中，我们引入\emph{Iwahori–Hecke 温性}这一概念。
它是关于 $\Aff_{G,x_0}$-模范畴的一个条件，蕴含张量函子
$\Whit^{!}(\Fl_G)_{x_0}\otimes_{\Aff_{G,x_0}}-$ 是保守的。

最后，在 §3.4.4 中，我们利用 $\B$ 与 $\IFLE$ 之间的相容性，
把 $\Whit^{!}(\Fl_G)_{x_0}\otimes_{\Aff_{G,x_0}}-$ 作用于 $\IFLE$ 的两侧，
并将结果嵌入环境范畴 $\IndCoh(\Op^{\mer}_{\Gd,x_0})$。
然后我们证明这两个嵌入彼此吻合，且 $\IFLE$ 的两侧都是 Iwahori–Hecke 温的，
从而完成 \thmref{thm:1.1.2.6} 的证明。

\srcpage{15}
